\par\SourceLeaf{421}%
\noindent\hspace*{1.2em}Thus, as one proceeds along the curve, the side toward which it is convex, or the signs in the formulas above, will remain unchanged as long as
\[
\frac{\delta V}{\delta\varrho}
\]
does not change sign. Insofar as $V$ is a continuous function, such a change of sign can occur only after passing through the value $0$. This necessarily requires that we have $p$ and $q=0$ simultaneously. At such a point the radius of curvature becomes infinitely large, or the curvature vanishes. Since, generally speaking,
\[
-p\sin\varphi+q\cos\varphi
\]
will change sign here, there is a point of inflection.

\Needspace{5\baselineskip}
\begin{center}5.\end{center}
\noindent\hspace*{1.2em}The case in which the nature of the curve is expressed by setting $y$ equal to a given function of $x$, namely $y=X$, is included in the preceding case if we set
\[
V=X-y
\]
If we put
\[
dX=X'dx,\qquad dX'=X''dx,
\]
we obtain
\[
\begin{gathered}
p=X',\qquad q=-1,\\
P=X'',\qquad Q=0,\qquad R=0,
\end{gathered}
\]
hence
\[
\pm\frac1r=\frac{X''}{(1+X'X')^{\frac32}}.
\]
Since $q$ is negative here, the upper sign applies as $x$ increases. We can therefore put it briefly as follows. With a positive $X''$ the curve is concave toward the same side as that on which the $y$-axis lies relative to the $x$-axis. With a negative $X''$ the curve is convex toward that side.
